\section{Discussion}\label{sec:discussion}

% Counts are taken from evaluation/d5_root_causes.csv, evaluation/analysis/report.md and robust_report.md.

\subsection{Hardening that breaks the workload it hardens}\label{sec:disc-hardening}

Under the production prompt, a recurring cause of failure is a security control applied without regard to the image it constrains. Eight configurations (three each from Claude and DeepSeek and two from GPT, all under the production prompt) hardened their database so that it could not start. In seven, the database container (PostgreSQL in six, MySQL in one) keeps its root-based entrypoint but has every Linux capability dropped: the entrypoint must take ownership of its data directory and then switch to an unprivileged user, and without the corresponding capabilities it can do neither. In the eighth, the database runs as a fixed non-root user with its data directory at the root of a volume it does not own, so initialisation cannot set the directory's permissions. These eight configurations score as well on Kubescape as the other production-prompt configurations (median compliance 89.6\% in both groups), so the scanner gives no sign of the defect.

Two further mechanisms follow the same pattern, each a recognised practice in isolation. Liveness or startup probes whose time budget is shorter than a slow-starting Java service needs kill each container before it finishes starting, so the services restart indefinitely; this accounts for six failures (five from Claude, four under the production prompt), observed directly in two (kubelet kills with no out-of-memory reason, and log lines showing the services still initialising) and inferred for the other four from the same probe settings and readiness record. A least-privilege NetworkPolicy that admits traffic to the Kubernetes Dashboard's metrics scraper only from the dashboard's own pods blocks the dashboard's default route to the scraper, which passes through the API server's service proxy; samples from the same model with an identical policy but a direct scraper address connected normally. This accounts for two failures, one supported by that matched contrast and the other inferred from the same policy, rather than by a controlled experiment. The third NetworkPolicy failure in Table~\ref{tab:root-causes} is a different mechanism: an Online Boutique policy admits internal traffic only from pods carrying a label that no pod carries.

The production prompt therefore buys static compliance partly at the cost of deployability (Section~\ref{sec:res-rq5}); static gains matter only for configurations that still deploy.

\subsection{Translating the Compose file too literally}\label{sec:disc-compose}

Every model receives the project's Docker Compose file, and models tend to translate it service by service. Where the Compose file reflects a Docker-host setup that the maintainers deliberately left out of their Kubernetes manifests, this faithfulness is a defect. Sock Shop's Compose file mounts the host's Docker socket into one service and runs a one-shot load generator; its Kubernetes manifest contains neither. Under the neutral prompt, 12 of the 15 Sock Shop configurations mount the socket (all three samples from Claude, DeepSeek and Gemini), against 4 of 15 under the production prompt, and 27 of the 30 include the load generator. Access to the Docker socket is equivalent to root on the node \citep{rahman2023k8smisconfig}, and the socket does not exist on nodes whose container runtime is not Docker Engine, as on the evaluation cluster. A load generator run as a long-lived Deployment restarts indefinitely, and in several samples its command invokes a start-up script that does not exist in the image. The load generator is the most frequent single cause of a configuration that serves correctly but is not fully Ready (Table~\ref{tab:root-causes}).

\subsection{What static analysis misses}\label{sec:disc-static}

Several failures were not visible to the static scanners. A memory limit too tight for a Java service is recorded by the cluster as an out-of-memory kill, but no scanner relates a limit to the workload it constrains. A database that cannot start under its security context passes every structural check, and so does a probe aimed at a path the image does not serve. Other defects are caught before deployment but are easy to miss in review: in two DeepSeek configurations, an unquoted RBAC resource name containing colons was parsed by YAML as a mapping, the API server rejected the whole Role, and the dashboard started without its permissions. Consistent with these cases, higher static security and reliability scores show no positive association with deployment success (Section~\ref{sec:res-rq4}); only requirement satisfaction, which checks that the documented components and settings are present, is weakly associated with it. Each of these defects surfaced only at deployment.

\subsection{``Ready'' is not the same as ``working''}\label{sec:disc-ready}

Kubernetes considers a workload Ready when its readiness probe passes, or immediately when it has none. The human Sock Shop manifest defines no readiness probes, so its slow Java services are reported Ready long before they accept connections, while several generated configurations serve every request correctly with an auxiliary workload not Ready. In this study the two outcomes disagree for 28 of the 198 generated configurations (18 serve correctly without being fully Ready and 10 are fully Ready without serving correctly). Reporting both readiness and a smoke test avoids rewarding missing probes and penalising non-essential components.

\subsection{Memorisation}\label{sec:disc-memorisation}

Two observations bear on the contamination threat (Section~\ref{sec:threats}). First, Gemini's provider withheld three responses for the Kubernetes Dashboard (two under the neutral prompt, one under the production prompt) because they reproduced training data too closely; the project's official manifest is widely published. Second, some models' configurations are more similar to the published human manifest than to the other models' configurations for the same project, which is the pattern expected if a published manifest is being reproduced. The excess is clearly positive for Claude on the Kubernetes Dashboard (Jaccard $+0.35$) and Emojivoto ($+0.23$) and for Gemini on Sock Shop ($+0.23$), and close to zero or negative elsewhere. For Claude, access to the human manifests during generation is ruled out by the session transcripts (Section~\ref{sec:method-generation}), so this similarity reflects what the model learned in training. Neither observation affects the scoring, but both limit how far results on widely published projects can be read as reasoning from documentation.

\subsection{Generalisation beyond the evaluated projects}\label{sec:disc-generalisation}

The failure mechanisms identified above arise from how generated configurations interact with common images and Kubernetes defaults rather than from any single application: database images with root-based entrypoints, slow-starting JVM services and Compose files written for a Docker host are widespread. We therefore expect these mechanisms to recur in similar projects. This expectation remains to be tested, and the per-project rates reported here should not be extrapolated to applications of other kinds.

\subsection{Implications}\label{sec:disc-implications}

For practitioners, the results suggest three checks before adopting a generated configuration: deploy it and exercise it, rather than relying on how it reads or scores; review every hardening setting against the documented behaviour of the image it constrains, in particular capability drops and fixed user IDs on database images and liveness probes on slow-starting services; and treat host-level settings carried over from a Compose file (socket mounts, host paths, one-shot services) as suspect. Asking a model for ``production-ready'' output makes these checks more necessary, not less. For tool builders, the results suggest scanner rules that relate a security setting to the image it is applied to. For researchers, benchmarks of generated infrastructure code should include execution, since a scanner score did not indicate which configurations would work.
